\begin{abstract}
LLM Agents route each task through model, tool and Web services, and their encrypted traffic leaks queries, task types and identities. Defenses for this traffic are emerging, but they are judged with website-fingerprinting methods that assume one link is the unit of protection and that cover traffic has no side effects. We systematize fifteen Agent studies (forty-four settings) against twenty-two traffic-analysis anchors and show that neither assumption holds. Tool feedback re-enters model requests, so a task is exposed on several independently controlled links. Agent cover is executed as application work, so its privacy result is incomplete without the task's outcome; no mapped privacy setting reports both an adapted attacker and matched utility. We then apply the resulting evaluation procedure to a controlled case, executed cover in an MCP Agent with two LLM planners, and it exposes a third condition. Cover inserted outside the model loop lacks the planner's decision latency: a label-free gap rule strips it and restores task recognition from 65--82 under the best size-only attack to 107--116 of 120 tasks. Drawing decoy timing from the planner's own history removes most of that advantage, but not all. Decoys remain separable through their destination schedule, and an attack that decodes removal jointly with the task still exceeds the best size-only attack in nine of ten cells, significantly in two. Cover must therefore imitate the planner's joint sequence and timing, not its timing marginals. We contribute access-aware attack and defense maps, a classical-to-Agent transfer matrix, and an evaluation procedure that binds residual inference to task, state and latency outcomes.
\end{abstract}
